\section{Calibration and representation refinement}
\label{sec:calibration}

\subsection{Independent episodes, dependent observations within episodes}

An \emph{episode} is one complete observed trajectory, or another complete
sampling unit specified by the application. Let $E_1,\ldots,E_n$ be
independent and identically distributed episodes, also independent of the
data used to construct and freeze the models. For horizon $h\in\mathcal H$,
let $\mathcal T_h$ be a fixed, nonempty set of causal forecast origins,
available in every episode and used by every compared candidate. Within an
episode, observations may be arbitrarily dependent and their distribution
may change with time.

Index the $J$ candidates by $j\in\{1,\ldots,J\}$ and write
$H=|\mathcal H|$ for the number of horizons. Candidate $j$ has a
representation $\Gamma_j$, a representation pseudometric $d_j$, and a
fixed forecaster $f_{jh}=g_{jh}\circ\Gamma_j$ at horizon $h$. Its decoder
$g_{jh}$ has a known Lipschitz budget $\kappa_{jh}$. Let $X_{i,t}$ denote
the history at origin $t$ in episode $i$, and let $y_{i,t+h}$ denote the
corresponding single-horizon observation (distinct from the joint target
$Y_t$ in Section~\ref{sec:method}). Define
\begin{equation}
 Z_{i,jh}=\frac1{|\mathcal T_h|}\sum_{t\in\mathcal T_h}
       \|y_{i,t+h}-f_{jh}(X_{i,t})\|_2^2,
 \qquad R_{jh}=\mathbb E Z_{i,jh}.
 \label{eq:episode-risk}
\end{equation}
Thus $Z_{i,jh}$ is one average loss per episode, and $R_{jh}$ is its
population mean. Equivalently, $R_{jh}$ is the risk for a new episode and
an independently, uniformly sampled origin from $\mathcal T_h$. It is not
a guarantee conditional on an individual history. Changing the episode
weights or the origin distribution changes the risk being certified.

Assume both targets and predictions lie almost surely in a known nonempty
closed convex set $\mathcal Y_h$. Let
$D_h=\sup_{y,z\in\mathcal Y_h}\|y-z\|_2<\infty$ be its diameter, and
put $B_h=D_h^2$. This gives the deterministic loss bound
$0\leq Z_{i,jh}\leq B_h$. Prediction clipping alone does not establish
this bound when the target is unbounded. Nor do training minima and maxima
establish a population range.

Form $K=\lfloor n/2\rfloor$ disjoint episode pairs, leaving one episode
unused for pairing when $n$ is odd. For $r\in\{1,\ldots,K\}$ define
\begin{equation}
\begin{split}
 Q_{r,jh}=\frac1{2|\mathcal T_h|}\sum_{t\in\mathcal T_h}
 \Bigl(&\|y_{2r-1,t+h}-y_{2r,t+h}\|_2\\
 &-\kappa_{jh}d_j\bigl(\Gamma_j(X_{2r-1,t}),
                         \Gamma_j(X_{2r,t})\bigr)\Bigr)_+^2.
\end{split}
 \label{eq:episode-pairs}
\end{equation}
For each $j,h$, the variables $Z_{i,jh}$ are independent across episodes,
and $Q_{r,jh}$ are independent across disjoint pairs. Terms within an
individual average need not be independent. Their ranges are $[0,B_h]$
and $[0,B_h/2]$, respectively. Write
$W_{jh}=\mathbb E Q_{r,jh}$. This is exactly the witness in
\eqref{eq:pair-witness} for the horizon-$h$ task when $A$ is a uniformly
sampled common origin: conditional on that origin, the two histories and
targets come from independent episodes.

For $N\geq2$ independent, identically distributed scalar observations
$V=(V_1,\ldots,V_N)$ in a known interval of length $b\geq0$, define the
sample mean and unbiased sample variance by
\[
 \bar V=\frac1N\sum_{i=1}^N V_i,
 \qquad s_V^2=\frac1{N-1}\sum_{i=1}^N(V_i-\bar V)^2.
\]
For a failure probability $\alpha\in(0,1)$, use the
empirical Bernstein radius
\begin{equation}
 \operatorname{rad}_{\rm EB}(V;b,\alpha)
 =\sqrt{\frac{2s_V^2\log(2/\alpha)}{N}}
       +\frac{7b\log(2/\alpha)}{3(N-1)}.
 \label{eq:eb-radius}
\end{equation}
Here $\log$ is the natural logarithm. The bounded-variable inequality of
\citet[Theorem~4]{maurer2009empirical}, also applied to the reflected observations
$-V_i$, gives each of the one-sided bounds
$\mathbb EV_1\leq\bar V+\operatorname{rad}_{\rm EB}$ and
$\mathbb EV_1\geq\bar V-\operatorname{rad}_{\rm EB}$ with failure
probability at most $\alpha$. A joint claim for both directions requires
allocation to both events.

\begin{theorem}[Simultaneous two-sided episode certificate]
\label{thm:episode-certificate}
Suppose the $J$ candidates and $H$ horizons are frozen using independent
development data, $n\geq4$, and all fitted decoders satisfy their stated
Lipschitz budgets. Fix $\delta\in(0,1)$ and set
$\alpha=\delta/(2JH)$. With
$\bar Q_{jh}=K^{-1}\sum_{r=1}^KQ_{r,jh}$ and
$\bar Z_{jh}=n^{-1}\sum_{i=1}^nZ_{i,jh}$, define
\begin{align}
 L_{jh}&=\max\{0,\bar Q_{jh}
 -\operatorname{rad}_{\rm EB}(Q_{\cdot,jh};B_h/2,\alpha)\},
 \label{eq:pair-lcb}\\
 U_{jh}&=\min\{B_h,\bar Z_{jh}
 +\operatorname{rad}_{\rm EB}(Z_{\cdot,jh};B_h,\alpha)\}.
 \label{eq:risk-ucb}
\end{align}
The dot denotes the full sample over episodes or pairs, so the radius uses
sample size $K$ for $Q_{\cdot,jh}$ and $n$ for $Z_{\cdot,jh}$. With
probability at least $1-\delta$, conditional on the development data,
all candidate--horizon combinations simultaneously satisfy
\begin{equation}
 L_{jh}\leq W_{jh}\leq R^*_{jh}\leq R_{jh}\leq U_{jh}.
 \label{eq:risk-sandwich}
\end{equation}
Here $R^*_{jh}$ is the infimum risk over all
$\kappa_{jh}$-Lipschitz decoders of $\Gamma_j$, using $d_j$, the same
horizon, the same uniform origin law, and the same raw target.
\end{theorem}

\begin{proof}
Condition on the development data so that the entire candidate family is
fixed. For each $j,h$, apply the upper empirical Bernstein bound to the
$n$ episode losses and the lower bound to the $K$ pair scores. Since
$0\leq W_{jh}$ and $R_{jh}\leq B_h$, truncating the endpoints at zero
and $B_h$ preserves validity. A union bound over these $2JH$ one-sided
events gives $L_{jh}\leq W_{jh}$ and $R_{jh}\leq U_{jh}$ jointly,
with failure probability at most $2JH\alpha=\delta$.

For any admissible decoder, Theorem~\ref{thm:pair-obstruction} applies to
the joint law obtained by drawing a common uniform origin and two
independent episodes. It bounds $W_{jh}$ by that decoder's risk under
the same origin mixture. Taking the infimum over these decoders gives
$W_{jh}\leq R^*_{jh}$. Finally, the fitted decoder belongs to this
class, so $R^*_{jh}\leq R_{jh}$. The loss and pair statistics may use
the same episodes: independence between the upper and lower statistics
is unnecessary for the union bound.
\end{proof}

The probability $1-\delta$ concerns repeated draws of the calibration
episodes. The theorem bounds expected squared error; it does not bound
every future error by $\sqrt{U_{jh}}$ or provide a prediction interval
with conditional coverage. Its independent sample sizes are $n$ and $K$,
not the number of overlapping windows. Additional origins may reduce the
variance of an episode average, but they do not create additional
independent episodes. The representation, normalization, dictionary sizes,
budgets, horizons, and all baselines included in the joint guarantee must
be fixed before calibration. They may depend on the development data
because the theorem conditions on those data.

\paragraph{Stratification under the same target law.}
A partition can concentrate pairs on distinctions that unconditional pairing
rarely encounters. Condition on the development data throughout this
construction, so the representation, metric, and budget are fixed.
Fix a measurable map
$c:\mathcal X\to\{1,\ldots,M\}$, with $M\geq1$, using development
data only, and write
$\pi_a=\Pr\{c(X)=a\}$. For each positive-mass stratum, let $W_a$ be
the witness from two independent draws conditional on $c(X)=a$.
Set $W_a=0$ when $\pi_a=0$. The stratified witness is
\begin{equation}
 W_{\rm str}=\sum_{a=1}^M\pi_a W_a\leq R^*_\kappa.
 \label{eq:stratified-witness}
\end{equation}
This is Theorem~\ref{thm:pair-obstruction} with $A=c(X)$; weighting by
the population masses preserves the original target law.

To estimate it, use $n\geq1$ independent, identically distributed observations
$(X_i,Y_i)$ from that law.
For each independent trajectory, draw one origin independently from the
prespecified common origin law. This produces independent observations under
the episode--origin law. Pooling all
dependent windows does not. The fitted-risk upper bound may still average
all prescribed origins within each episode; a union bound does not require
independence between that bound and the stratified statistic.
Let $N_a=\sum_i\mathbf1\{c(X_i)=a\}$, where $\mathbf1\{\cdot\}$
denotes an indicator, and set $K_a=\lfloor N_a/2\rfloor$.
Set $U_i=\Gamma(X_i)$ and write $(U_{a,k},Y_{a,k})$,
$k=1,\ldots,N_a$, for the observations in stratum $a$, listed in their
original sample-index order. For $r=1,\ldots,K_a$, their disjoint pair
scores are
\[
 Q_{a,r}=\frac12\left(\|Y_{a,2r-1}-Y_{a,2r}\|_2
      -\kappa d(U_{a,2r-1},U_{a,2r})\right)_+^2.
\]
Assume the target set has squared diameter $B<\infty$, so
$0\leq Q_{a,r}\leq B/2$.

\begin{theorem}[A stratified lower certificate with estimated masses]
\label{thm:stratified-certificate}
Fix $\delta_p,\beta_1,\ldots,\beta_M\in(0,1)$ with
$\delta_p+\sum_a\beta_a<1$. If $K_a<2$, set $\ell_a=0$; otherwise set
\[
 \ell_a=\max\{0,\bar Q_a-
       \operatorname{rad}_{\rm EB}(Q_{a,\cdot};B/2,\beta_a)\}.
\]
Here $\bar Q_a$ is the mean of the $K_a$ scores. Define
\[
 t_p=\sqrt{\frac{\log(2M/\delta_p)}{2n}},\qquad
 p_a^-=(N_a/n-t_p)_+,\qquad
 p_a^+=\min\{1,N_a/n+t_p\},
\]
and the confidence set, with
$\Delta_M=\{p\in[0,1]^M:\sum_{a=1}^M p_a=1\}$ the probability simplex,
$\mathcal P=\{p\in\Delta_M:p_a^-\leq p_a\leq p_a^+
\text{ for every }a\}$. Then
\begin{equation}
 L_{\rm str}:=\min_{p\in\mathcal P}\sum_a p_a\ell_a
 \ \leq\ W_{\rm str}\ \leq\ R^*_\kappa
 \label{eq:stratified-certificate}
\end{equation}
with probability at least $1-\delta_p-\sum_a\beta_a$.
\end{theorem}

\begin{proof}
Condition on the complete stratum-label vector. The observations in each
stratum are independent draws from its conditional law; their disjoint pair
scores are therefore independent and identically distributed. For every
realized count with $K_a\geq2$, the conditional failure probability of
$\ell_a\leq W_a$ is at most $\beta_a$. Sparse strata satisfy it
deterministically. Averaging over the labels and taking a union bound gives
all score inequalities with failure probability at most $\sum_a\beta_a$;
there is no additional union over possible counts.
For each stratum, Hoeffding's inequality gives
$\Pr\{|N_a/n-\pi_a|>t_p\}\leq2\exp(-2nt_p^2)=\delta_p/M$.
A union bound over the $M$ strata therefore places
$\pi=(\pi_1,\ldots,\pi_M)$ in $\mathcal P$ with failure probability
at most $\delta_p$.
On the intersection of these events,
$\min_{p\in\mathcal P}\sum_a p_a\ell_a
\leq\sum_a\pi_a\ell_a\leq\sum_a\pi_a W_a$.
No independence between masses and score bounds is used.
\end{proof}

The set $\mathcal P$ is nonempty because it contains $(N_a/n)_a$.
Its linear minimum is computed by starting at $p^-$ and assigning the
remaining mass, up to each upper bound, in increasing order of $\ell_a$.
Known population masses instead give $L_{\rm str}=\sum_a\pi_a\ell_a$
with no mass-estimation allocation. For several frozen partitions or
candidate--horizon combinations, allocate failure probability over their
score events and mass events before taking the largest lower endpoint.
A shared partition needs only one mass event. Sparse strata contribute no
positive evidence; stratification does not increase the number of independent
observations.

\begin{proposition}[Exactness on compatible strata]
\label{prop:stratified-exactness}
Assume $\mathbb E\|Y\|_2^2<\infty$. Suppose each positive-mass
stratum either has a constant representation
$U=u_a$ or a joint law supported on two values of $(U,Y)$.
In the first case let its optimal prediction be
$\mu_a=\mathbb E[Y\mid c(X)=a]$; in the second use the fitted values
from Proposition~\ref{prop:two-atom}, with $b=d(u_0,u_1)$ in that
construction. If a single global
$\kappa$-Lipschitz decoder realizes all these values, then
$W_{\rm str}=R^*_\kappa$. For scalar targets on finite support, global
compatibility is equivalent to the pairwise Lipschitz inequalities between
all the proposed fitted values, including values from different strata.
\end{proposition}

\begin{proof}
A constant representation gives
$W_a=\mathbb E[\|Y-\mu_a\|_2^2\mid c(X)=a]$ by the independent-pair
variance identity. A two-atom stratum has the finite-support calculation of
Proposition~\ref{prop:two-atom}. The assumed decoder attains their weighted
sum, giving equality in \eqref{eq:stratified-witness}. For scalar prescribed
values $v_i$ at finitely many support points $u_i$, define
$g(u)=\min_i\{v_i+\kappa d(u,u_i)\}$. Each function in this finite
minimum is $\kappa$-Lipschitz by the triangle inequality for $d$, and
so is their minimum. At $u_j$, the term $i=j$ equals $v_j$ and every
other term is at least $v_j$ by the assumed pairwise inequalities.
Thus $g(u_j)=v_j$. Conversely, a global Lipschitz decoder necessarily
satisfies all these pairwise inequalities.
\end{proof}

For the three-atom example, the stratum $\{3\}$ has zero witness.
The stratum $\{0,1\}$ has mass $2\varepsilon$, and its two atoms
are conditionally equally likely, giving conditional witness $1/4$.
Consequently, $W_{\rm str}=2\varepsilon(1/4)=\varepsilon/2=R^*_\kappa$.
Stratification can thus remove a population gap, rather than merely improve
concentration. It need not improve every witness: singleton strata can remove
all pairwise conflicts while incompatible predictions across strata still
force positive global risk. Conditional optima cannot be combined without
the global compatibility condition.

\begin{corollary}[Acceptance, exclusion, and inconclusive decisions]
\label{cor:decisions}
Let $\epsilon_h\geq0$ be a prespecified root-mean-square error (RMSE)
tolerance at horizon $h$. On the event in \eqref{eq:risk-sandwich},
$U_{jh}\leq\epsilon_h^2$ certifies that the fitted forecaster attains
this tolerance. If $L_{jh}>\epsilon_h^2$, no decoder in the corresponding
Lipschitz class attains it. When neither condition holds, the decision is
inconclusive.
Selecting any member of the frozen family that satisfies
$U_{jh}\leq\epsilon_h^2$ at every required horizon preserves these
guarantees. In particular, one may minimize a predeclared cost over the
certified candidates and abstain if there are none.
\end{corollary}

\begin{corollary}[A certificate of minimum feasible cost]
\label{cor:cost-optimality}
Let candidate $j$ represent a class $\mathcal C_j$ of multihorizon forecasts
through $\Gamma_j$, containing its fitted forecast and satisfying its stated horizon-specific
Lipschitz budgets. Give it a fixed cost $c_j\geq0$. Let $c^*$ be the minimum
cost of a class containing a forecast that meets every prescribed tolerance.
The minimum of an empty set is $+\infty$.
For any simultaneously valid lower bounds $L_{jh}$ on the horizon-$h$
risk of every member of $\mathcal C_j$, and upper bounds $U_{jh}$ on
its fitted forecast, define
\begin{align*}
 \mathcal E&=\{j:L_{jh}>\epsilon_h^2\text{ for some }h\},&
 \mathcal A&=\{j:U_{jh}\leq\epsilon_h^2\text{ for all }h\},\\
 c_-&=\min_{j\notin\mathcal E}c_j,&
 c_+&=\min_{j\in\mathcal A}c_j.
\end{align*}
Thus $\mathcal E$ contains the excluded classes and $\mathcal A$ contains
the accepted fitted candidates. The costs $c_-$ and $c_+$ refer to the
cheapest nonexcluded class and cheapest accepted fit, respectively.
On their joint confidence event,
\begin{equation}
 c_-\leq c^*\leq c_+.
 \label{eq:cost-interval}
\end{equation}
If $c_-=c_+<\infty$, a least-cost accepted forecast has minimum feasible
cost among the specified classes. If every class is excluded, the tolerance
requirements are infeasible within their union.
\end{corollary}

\begin{proof}
Every accepted fitted forecast meets the tolerances. Every class containing
a forecast that meets them is nonexcluded, since one violated horizon lower
bound would rule out that forecast. These two inclusions give
\eqref{eq:cost-interval} and both conclusions. A joint coefficient budget
can further restrict $\mathcal C_j$; exclusion by its larger horizon-wise
Lipschitz classes remains valid.
\end{proof}

If $c_-<c_+$, a cheaper nonexcluded class may still contain an accurate
forecast, so selection establishes only the lowest cost among certified
candidates. In the experiments, cost counts feature coordinates rather
than runtime.
Also, \eqref{eq:risk-sandwich} implies
\[
 0\leq\sqrt{R_{jh}}-\sqrt{R^*_{jh}}
       \leq\sqrt{U_{jh}}-\sqrt{L_{jh}}.
\]
Thus the endpoint gap bounds the fitted decoder's RMSE suboptimality within
the stated class. A large gap does not identify optimization error: the
witness and either confidence correction may be loose.

For simple sample-size calculations, one may instead prespecify the
Hoeffding radius $b\sqrt{\log(1/\alpha)/(2N)}$. If a fitted risk lies
at least $\gamma>0$ below a squared-error tolerance, then
\[
 n\geq\frac{2B_h^2\log(1/\alpha)}{\gamma^2}
\]
is sufficient for its Hoeffding upper bound to pass that tolerance with
probability at least $1-\alpha$. If the pair witness exceeds the tolerance
by $\gamma$, then
\[
 K>\frac{B_h^2\log(1/\alpha)}{2\gamma^2}
\]
is sufficient for the corresponding lower bound to exclude the tolerance
with that probability. For acceptance, the stated sample size makes the
radius at most $\gamma/2$. A one-sided Hoeffding event bounds the sample
mean by the population mean plus $\gamma/2$, so the upper endpoint is
at most the tolerance. For exclusion, the strict sample-size inequality
makes the radius strictly smaller than $\gamma/2$. On the event $\bar Q\geq W-\gamma/2$, where $W$ is the population
pair witness and $\bar Q$ is its sample mean, the lower endpoint is
therefore strictly above the tolerance. Each event has probability at least $1-\alpha$.
These are sufficient sample sizes, not optimal rates. Exclusion depends on the witness margin, rather than on an unknown
oracle-risk margin. Choosing the smaller of separately computed Hoeffding
and empirical Bernstein bounds requires confidence allocation to both
bounds; otherwise the choice must be fixed in advance.

\subsection{Refinement that preserves the meaning of the norm budget}

Let $\Phi_b:\mathcal X\to\Delta_{C_b}$, $b=1,\ldots,s$, be fixed
dictionary blocks, where
$\Delta_{C_b}=\{u\in\mathbb R^{C_b}:u_c\geq0,\ \sum_cu_c=1\}$
is the probability simplex. Retain the entire stack
$\Gamma_s=(\Phi_1,\ldots,\Phi_s)$. Refinement appends a block while
keeping every previous coordinate function unchanged. For stacked
vectors $u=(u_1,\ldots,u_s)$ and $v=(v_1,\ldots,v_s)$ define the block
maximum metric $\rho_s(u,v)=\max_{1\leq b\leq s}\|u_b-v_b\|_2$.
Its induced pseudometric on histories and the coefficient norm are
\begin{equation}
 d_s(x,x')=\rho_s(\Gamma_s(x),\Gamma_s(x'))
          =\max_{1\leq b\leq s}\|\Phi_b(x)-\Phi_b(x')\|_2,
 \qquad
 \|W\|_{2,1}=\sum_{b=1}^s\|W_b\|_{\rm F}.
 \label{eq:block-geometry}
\end{equation}
Here $\|W\|_{2,1}$ is a sum over \emph{blocks}, not over individual
rows. As in \eqref{eq:training-objective}, $a\in\mathbb R^q$ and
$W_b\in\mathbb R^{C_b\times q}$, so the contribution of block $b$
to a column-vector prediction is $W_b^\top\Phi_b(x)$.

Let $\mathcal Y\subseteq\mathbb R^q$ be a nonempty closed convex output
set containing $Y$ almost surely. Write
$\Pi_{\mathcal Y}(z)=\arg\min_{y\in\mathcal Y}\|z-y\|_2$ for Euclidean
projection; for a box, it clips each coordinate to its allowed interval.
Define
\begin{equation}
 \mathcal F_{s,\kappa}
 =\left\{x\mapsto\Pi_{\mathcal Y}\!\left(
       a+\sum_{b=1}^sW_b^\top\Phi_b(x)\right):
       a\in\mathbb R^q,\ \sum_{b=1}^s\|W_b\|_{\rm F}\leq\kappa\right\}.
 \label{eq:block-decoder}
\end{equation}
The intercept is unconstrained. Group norm constraints are established
tools \citep{bach2008group}; here the block norm makes the decoder's
sensitivity compatible with retained dictionary blocks.

\begin{theorem}[Nested and replication-invariant refinement]
\label{thm:replication-invariance}
Every decoder in $\mathcal F_{s,\kappa}$ is $\kappa$-Lipschitz with
respect to $d_s$. Appending a block gives
$\mathcal F_{s,\kappa}\subseteq\mathcal F_{s+1,\kappa}$ and
$d_{s+1}\geq d_s$ pointwise. Consequently, the infimum prediction risk
in this decoder class cannot increase. Neither can the pair witness,
provided the budget, target law, and conditional pairing law remain fixed.
Appending an exact copy of an existing block leaves both the decoder
class and the pair witness unchanged.
\end{theorem}

\begin{proof}
For the unprojected decoder, the triangle inequality and the bound
$\|Av\|_2\leq\|A\|_{\rm F}\|v\|_2$ give
\[
 \left\|\sum_bW_b^\top\bigl(\Phi_b(x)-\Phi_b(x')\bigr)\right\|_2
 \leq\sum_b\|W_b\|_{\rm F}\|\Phi_b(x)-\Phi_b(x')\|_2
 \leq\kappa d_s(x,x').
\]
Projection onto a closed convex set is nonexpansive, so this bound also
holds for the projected decoder. Assigning zero coefficients to a new
block embeds each old decoder in the refined class, proving inclusion and
hence nonincrease of the optimal risk. Appending a block can only increase
the maximum in $d_s$. For each paired observation, increasing that distance
can only decrease the positive-part discrepancy in
\eqref{eq:pair-witness}; taking expectations proves witness monotonicity.

If block $s+1$ exactly copies block $b$, merge its coefficients into the
old block. The prediction is unchanged and
$\|W_b+W_{s+1}\|_{\rm F}\leq\|W_b\|_{\rm F}+\|W_{s+1}\|_{\rm F}$,
so the merged coefficients remain feasible. This proves the reverse
class inclusion for duplication. A duplicate also leaves the maximum
distance unchanged, and therefore leaves the witness unchanged.
\end{proof}

These claims concern population optima over nested feasible classes, not
the risks of separately fitted finite-sample estimators. They also do not
assert that every new block contains useful predictive information.
Since the Euclidean diameter of each probability simplex is at most $\sqrt2$,
$d_s\leq\sqrt2$ regardless of the number of retained blocks. Certifying a
fitted decoder at its smaller actual norm is possible, but changes the
comparator class. Refinement monotonicity is asserted only at a common
budget.

The choice of norm matters even when no information is added. Let two
equally likely histories satisfy $\Phi(x_0)=e_1=(1,0)^\top$ and
$\Phi(x_1)=e_2=(0,1)^\top$, with scalar targets $-1$ and $1$. Under a
single Frobenius constraint of radius one, the largest possible forecast
separation is $\sqrt2$, and the best affine mean squared error (MSE) is
$(2-\sqrt2)^2/4$. After duplicating the block, the concatenated coefficient
vector $(-1/2,1/2,-1/2,1/2)$ has norm one and predicts both targets exactly.
Duplication has relaxed the effective sensitivity constraint. The
block-sum constraint prevents this effect while allowing genuinely new
blocks to change the geometry.

\begin{theorem}[A necessary amount of geometric change]
\label{thm:necessary-expansion}
Let $d_0,d_1$ be measurable pseudometrics induced on a common history space by an
old and a new representation. Use the same square-integrable target,
conditionally independent pairing law, and sensitivity budget $\kappa>0$.
Define the random distance increase
$\Delta=(d_1(X,X')-d_0(X,X'))_+$ and assume
$\|\Delta\|_{L^2}<\infty$. Let $W_0,W_1$ be the corresponding witnesses
from \eqref{eq:pair-witness}, using these induced history distances. Then
\begin{equation}
 \sqrt{W_1}\geq
 \left(\sqrt{W_0}-\frac{\kappa}{\sqrt2}
                       \|\Delta\|_{L^2}\right)_+.
 \label{eq:metric-change}
\end{equation}
In particular, if a $\kappa$-Lipschitz decoder of the new representation
attains RMSE at most $\epsilon\geq0$, then
\begin{equation}
 \|\Delta\|_{L^2}\geq
       \frac{\sqrt2}{\kappa}(\sqrt{W_0}-\epsilon)_+.
 \label{eq:required-change}
\end{equation}
On any event $0\leq L_0\leq W_0$, replacing $W_0$ by $L_0$ on the
right side of \eqref{eq:required-change} is valid simultaneously for every
new representation satisfying these conditions. The new representation
need not have been fixed when $L_0$ was computed.
\end{theorem}

\begin{proof}
Set $a=\|Y-Y'\|_2$ and
$v_i=(a-\kappa d_i(X,X'))_+$ for $i=0,1$. The positive-part function
is nondecreasing and $1$-Lipschitz. Consequently,
$v_0\leq v_1+\kappa\Delta$ pointwise, and the $L^2$ triangle inequality
gives
\[
 \|v_1\|_{L^2}\geq
       (\|v_0\|_{L^2}-\kappa\|\Delta\|_{L^2})_+.
\]
Since $\|v_i\|_{L^2}=\sqrt{2W_i}$, dividing by $\sqrt2$ proves
\eqref{eq:metric-change}. If the new decoder has RMSE at most $\epsilon$,
Theorem~\ref{thm:pair-obstruction} gives $W_1\leq\epsilon^2$.
Combining the two inequalities and rearranging gives
\eqref{eq:required-change}. Finally, its right side is nondecreasing in
$W_0$. Replacing $W_0$ by a nonnegative lower bound therefore weakens
the necessary condition, for every admissible new representation on the
same event.
\end{proof}

This theorem rules out a proposed repair when its population distance
increase is too small. It is a necessary condition, not a sufficient
refinement rule or a guarantee that an estimated distance increase is
accurate. The quantity $L_0$ must lower-bound the old \emph{pair witness};
an arbitrary lower bound on old oracle risk cannot be substituted.
If the budget also changes from $\kappa_0$ to $\kappa_1$, the same proof
gives the necessary condition
\[
 \bigl\|(\kappa_1d_1(X,X')-\kappa_0d_0(X,X'))_+\bigr\|_{L^2}
 \geq\sqrt2(\sqrt{W_{\kappa_0,0}}-\epsilon)_+,
\]
where $W_{\kappa_0,0}$ is the old witness at budget $\kappa_0$.
This form allows the increased separation to come from a larger budget,
a changed representation, or both.

\subsection{When an added block can improve prediction}

The preceding theorem gives a necessary geometric condition. The next
calculation gives a sufficient population condition for one proposed
feature block: its cross-moment with the current residual must be nonzero,
and the old decoder must have unused norm budget. A positive obstruction
alone implies neither condition.

\begin{proposition}[Residual correlation and budget slack]
\label{prop:residual-gain}
Assume $\mathbb E\|Y\|_2^2<\infty$ and consider the unprojected affine
predictor
\[
 f_0(x)=a+\sum_{b=1}^sW_b^\top\Phi_b(x),
 \qquad \sum_{b=1}^s\|W_b\|_{\rm F}<\kappa.
\]
Let $\psi:\mathcal X\to\mathbb R^p$ be a new feature block with
$\mathbb E\|\psi(X)\|_2^2<\infty$. This block may take general
Euclidean values; it need not be simplex-valued. The same sum of block
Frobenius norms constrains its coefficients. Define
\begin{align*}
 r&=Y-f_0(X),
 &s_0&=\kappa-\sum_b\|W_b\|_{\rm F}>0,\\
 C&=\mathbb E[\psi(X)r^\top]\in\mathbb R^{p\times q},
 &G&=\mathbb E[\psi(X)\psi(X)^\top]\in\mathbb R^{p\times p}.
\end{align*}
Here $C$ is the feature--residual cross-moment, $G$ is the feature second
moment matrix, and $s_0$ is the unused budget. If $C\ne0$, let
$a_0=\|C\|_{\rm F}^2$ and $b_0=\operatorname{tr}(C^\top GC)$, where
$\operatorname{tr}$ sums the diagonal entries. Then $b_0>0$, and the
coefficient block
\[
 V=\eta C\in\mathbb R^{p\times q},\qquad
 \eta=\min\{a_0/b_0,\ s_0/\sqrt{a_0}\}
\]
is feasible when appended to the old blocks. Its risk improvement is
\begin{equation}
 R(f_0)-R(f_0+V^\top\psi)
       =2\eta a_0-\eta^2b_0\geq\eta a_0>0.
 \label{eq:residual-gain}
\end{equation}
If $Y$ and $f_0(X)$ belong almost surely to a nonempty closed convex set
$\mathcal Y$, projecting the new forecast onto $\mathcal Y$ retains at
least this improvement.
\end{proposition}

\begin{proof}
The old simplex-valued blocks are bounded, so $f_0(X)$ and hence the
residual $r$ have finite second moments. The same assumption on the new
block makes $C$ and $G$ well-defined by the Cauchy--Schwarz inequality. Also,
$b_0=\mathbb E\|C^\top\psi(X)\|_2^2$. If $b_0=0$, then
$C^\top\psi(X)=0$ almost surely, implying
\[
 \|C\|_{\rm F}^2
 =\langle C,\mathbb E[\psi(X)r^\top]\rangle_{\rm F}
 =\mathbb E[r^\top C^\top\psi(X)]=0,
\]
contrary to $C\ne0$. Thus $b_0>0$ and $\eta>0$.

For any coefficient matrix $V\in\mathbb R^{p\times q}$, expand
$\|r-V^\top\psi(X)\|_2^2$ and take expectations to obtain
\[
 R(f_0+V^\top\psi)=R(f_0)-2\langle V,C\rangle_{\rm F}
                         +\operatorname{tr}(V^\top GV).
\]
Substituting $V=\eta C$ gives the equality in
\eqref{eq:residual-gain}. Since
$\|V\|_{\rm F}=\eta\sqrt{a_0}\leq s_0$, appending this block respects
the total budget. Since $\eta\leq a_0/b_0$, the improvement is at least
$\eta a_0>0$. Finally, projection onto a closed convex set containing
$Y$ cannot increase the new squared error. The old forecast is already
in that set, so its reference risk is unchanged.
\end{proof}

The condition on the old forecast matters. An improvement between two
unprojected affine predictors does not by itself imply an improvement
between their separately clipped versions. The proposition also makes no
claim when the old decoder exhausts its budget. Its moments are population
moments: replacing them with training estimates does not establish the
same population improvement. For fitted, clipped models, we therefore
compare their episode losses directly.

Fix a horizon $h$ and two forecasts before collecting $n\geq2$ new
independent evaluation episodes. Let $Z_i^{\rm old}$ and
$Z_i^{\rm new}$ be their losses on episode $i$, averaged over the same
origins as in \eqref{eq:episode-risk}. Define
\[
 D_i=Z_i^{\rm old}-Z_i^{\rm new}\in[-B_h,B_h],\qquad
 \bar D=\frac1n\sum_{i=1}^nD_i,\qquad
 R^{\rm old}=\mathbb EZ_i^{\rm old},\quad
 R^{\rm new}=\mathbb EZ_i^{\rm new}.
\]
For $\alpha\in(0,1)$, the lower empirical Bernstein bound gives, with
probability at least $1-\alpha$,
\begin{equation}
 R^{\rm old}-R^{\rm new}\geq
 \bar D-\operatorname{rad}_{\rm EB}(D;2B_h,\alpha).
 \label{eq:paired-improvement}
\end{equation}
The interval length is $2B_h$, because $D_i$ can be negative. A positive
right side certifies a risk reduction for these two fitted predictors;
a nonpositive value is inconclusive. For multiple fixed comparisons,
allocate failure probability over the entire comparison family.

Sequential refinement can also be adaptive. At stage $s$, a candidate may
depend on all previous results. Conditional on this full past, its evaluation
episodes must be new independent draws from the specified evaluation law.
Choose stage failure probabilities $\alpha_s>0$ in advance, with
$\sum_s\alpha_s\leq\delta$, and accept
a change only when its paired lower bound is positive. Conditional validity
at each stage and a union bound then make every accepted improvement valid
on one event of probability at least $1-\delta$. This argument does not
permit repeated tuning on a fixed calibration set.

Finally, a positive obstruction need not be removable by refinement. If
$Y$ is independent of $X$ and takes values $-a$ and $a$ with equal
probability, where $a>0$, then the minimum risk over all predictors based
on $X$ is $a^2$. Under unconditioned independent pairing, a constant
representation has pair witness $a^2$; no representation of the same
observed history removes this independent future noise. Conversely, a zero
witness need not imply a useful predictor. For example, suppose every
episode follows the same deterministic trajectory,
but its target changes across forecast origins. Pairing two episodes at the
same origin gives identical targets and hence a zero witness. A constant
representation still cannot predict the variation across origins; its optimal
risk is $\mathbb E\|Y-\mathbb EY\|_2^2$ under the origin law.
Obstruction, fitted risk, and inconclusive decisions therefore answer
distinct questions and should be reported separately.
